\section{C. Congruence quotients and arithmetic comparisons}

\begin{claim}{C.r1}{Proved}{Exact arithmetic outer images on every congruence quotient}
For $n\ge2$, put $R_n=\mathbb Z/n$ and $K_n=\SL_2(R_n)$.
The outer action of $G_\mathbb Q$ on $K_n$ induced by the universal
elliptic curve is the cyclotomic character followed by an injection
\[
 R_n^\times/(R_n^\times)^2\hookrightarrow\Out(K_n).
 \tag{C.r1.1}
\]
Its image is exactly this square-class group. At the adelic level its
kernel is $\chi^{-1}((\Zhat^\times)^2)$, the subgroup fixing the
maximal multiquadratic extension of $\mathbb Q$. This proves the
arithmetic statements of \S39, scans22--24, folios299--300, with
``outer'' retained throughout.
\end{claim}
\begin{proof}
Over $\mathcal M_{1,1}$ the full $n$-torsion local system gives a
continuous representation
$\pi_1(\mathcal M_{1,1})\to\GL_2(R_n)$.
The Weil pairing identifies its determinant with $\chi\bmod n$;
over $\overline{\mathbb Q}$ the topological monodromy is
$\SL_2(\mathbb Z)\twoheadrightarrow K_n$.
Thus the exact arithmetic stack extension induces on $K_n$ the outer
action given by conjugation by a matrix of determinant $\chi(\gamma)$.
Different lifts of $\gamma$ change that matrix by an element of $K_n$.

A matrix commuting with the elementary matrices
$\left(\begin{smallmatrix}1&1\\0&1\end{smallmatrix}\right)$ and
$\left(\begin{smallmatrix}1&0\\1&1\end{smallmatrix}\right)$ is scalar:
entrywise comparison forces both off-diagonal entries to vanish and
the diagonal entries to coincide. Therefore the kernel of
$\GL_2(R_n)\to\Out(K_n)$ is exactly
$K_n\cdot(R_n^\times I)$. Taking determinants identifies the quotient
with $R_n^\times/(R_n^\times)^2$ and proves injectivity in (C.r1.1).
Surjectivity of the cyclotomic character modulo $n$ follows from
$\Gal(\mathbb Q(\zeta_n)/\mathbb Q)=R_n^\times$, so the stated outer
image is exact.

The same centralizer proof applies over $\Zhat$. Its unit group is a
product of compact groups; the subgroup of squares is the product of
their square subgroups and is closed. By the cyclotomic description of
the abelianized absolute Galois group (Kronecker--Weber), its quotient
by squares is the maximal elementary abelian two quotient. Equivalently
its fixed field is the compositum of all quadratic fields. This last
identification can also be read characterwise from the quadratic
Dirichlet characters. This proves all the asserted kernel and image
statements; it does not imply that the action on the full profinite
modular group factors through this small quotient.
\end{proof}



\begin{center}\small
\begin{tabular}{rrrrrr}
\toprule
$n$&$|\SL_2(R_n)|$&$|Z(\SL_2(R_n))|$&$|\Ad\GL_2(R_n)|$
 &$|\Inn(K_n)|$&arithmetic outer image\\
\midrule
2&6&1&6&6&1\\
3&24&2&24&12&$C_2$\\
4&48&2&48&24&$C_2$\\
5&120&2&120&60&$C_2$\\
6&144&2&144&72&$C_2$\\
7&336&2&336&168&$C_2$\\
8&384&4&384&96&$C_2^2$\\
12&1152&4&1152&288&$C_2^2$\\
24&9216&8&9216&1152&$C_2^3$\\
\bottomrule
\end{tabular}\end{center}
All entries, every unit square, and the scalar-centralizer assertion are
exhaustively checked in \nolinkurl{finite_congruence.py}. Notice the
centers of orders four and eight: replacing them by a uniform
$\{\pm I\}$ would give wrong inner-automorphism counts.

\begin{claim}{C.r2}{Refuted}{A unique multiplier cannot be recovered from an outer congruence action}
The claim on \S39, scans25--26, folios300 verso--301, that a unique
$\chi(u)\in\Zhat^\times$ is recovered even from an outer automorphism
of $\SL_2(\Zhat)$ preserving its unipotent cusp subgroup is false.
The strongest conclusion justified by this datum is its square class,
as in C.r1. On the full profinite modular group, the peripheral
multiplier is well defined; the passage to the congruence completion
destroys precisely the indicated uniqueness argument.
\end{claim}
\begin{proof}
For a unit $a$, conjugation by
$D_a=\operatorname{diag}(a,a^{-1})\in\SL_2$ sends
$U(t)=\left(\begin{smallmatrix}1&t\\0&1\end{smallmatrix}\right)$ to
$U(a^2t)$. It is an inner automorphism, so represents the same outer
class as the identity, while its multiplier on the displayed cusp
parametrization is $a^2$. Modulo five, take $a=2$; the two multipliers
are $1$ and $4$. Adelic examples are obtained by taking $a_5=2$ and
$a_p=1$ at the other primes. This is an explicit counterexample to
uniqueness, already in the requested finite range.

For comparison, in the full projective modular group let
$T=sr$, whose square belongs to the index-six free group $F$.
An inner conjugator normalizing $\langle T\rangle$ has reduction in
the stabilizer of the corresponding cusp in $S_3$. That stabilizer is
generated by the reduction of $T$, so it may be removed by a power
of $T$. The remaining conjugator is in $F$. Abelianizing $F$ then
shows that a conjugacy $T^2\mapsto (T^2)^b$ forces $b=1$, because
the cusp vector in $F^{\rm ab}\simeq\Zhat^2$ is primitive. Thus this
ambiguity does not occur in the original profinite peripheral structure.
The same argument applies to its linear central lift, whose extra
central subgroup acts trivially by conjugation.
\end{proof}

\begin{claim}{C.r3}{Proved}{Genuine and outer arithmetic actions must be counted separately}
The input $G_\mathbb Q\to\Out(\widehat{\SL_2(\mathbb Z)})$ has a
canonical outer meaning. Without choosing a section and compatible
paths it does not specify a homomorphism to the genuine automorphism
group. At a finite congruence level, the full arithmetic fundamental
group acts with image $\Ad\GL_2(R_n)$, the third numeric column above;
Galois acts canonically with the square-class outer image in the last
column. These are different requested objects, not alternative counts
for the same action.
\end{claim}
% anchor:C.r3
\manuscriptanchor{\S \S 39,45; fols.299--301,387--388; scans22--26,176--178}

\begin{proof}
The stack extension is an extension of $G_\mathbb Q$ by the geometric
modular group. Conjugation by a chosen lift of $\gamma$ is changed by
an inner automorphism if that lift is changed. This gives the outer
action and explains its ambiguity directly. Its torsion representation
contains every determinant-one matrix by geometric monodromy and has
every determinant by cyclotomic surjectivity; consequently its image
is all of $\GL_2(R_n)$. Quotient by its scalar centralizer gives the
genuine action image, and quotient further by inner conjugations gives
C.r1. The distinction survives at inverse limits. The faithful action
on the full profinite modular group and its centralizer are proved in
C.a5--C.a7, rather than inferred from finite congruence counts.
\end{proof}

\subsection{The congruence representation obtained from actual profinite cocycles}

\begin{claim}{C.d1}{Known}{Projective congruence kernels are preserved by $\widehat{\mathrm{GT}}_0$}
Let $q:\widehat B\twoheadrightarrow
\operatorname{SL}_2(\widehat{\mathbb Z})/\{\pm I\}$ be the congruence homomorphism, with the single global central subgroup $\{\pm I\}$. For every $u\in\widetilde M^!$, there is $U(u)\in\operatorname{GL}_2(\widehat{\mathbb Z})/\{\pm I\}$ of determinant $\mu(u)$ such that
\[
 q\circ u=\operatorname{Int}(U(u))\circ q.
\]
Consequently all projective congruence kernels are invariant under the genuine extension of $\widehat{\mathrm{GT}}_0$. This statement needs neither the pentagon nor an arithmetic-image assumption. It does not assert a canonical action on $\widehat{\operatorname{SL}_2(\mathbb Z)}$ without the specified central pullback. Its profinite structural input is the torsion theorem for free profinite products, not an inference from a matrix cocycle at one finite level.
\end{claim}
% anchor:C.d1
\manuscriptanchor{\S 39; scans21--30; \S 44; scans150--168; \S 46; fols398--400,409; scans198--202,218--219}

\begin{proof}
We make the structural input explicit. In $\widehat B=\widehat C_2\amalg\widehat C_3$, every finite-order element is conjugate into a factor (Herfort--Ribes, 1985, Theorem 2(i) \citep{HerfortRibes1985dictionary}, also explicitly used in Minamide--Nakamura, p.8, in the attached paper). The primary Herfort--Ribes proof was not supplied; this imported theorem is not a proof certified by the present calculation. Applying it to $\eta s$ and $\xi r$ in B.m1 gives
\[
 \eta=\beta\sigma(\beta)^{-1},\qquad
 \xi=\alpha\rho(\alpha)^{-1},\qquad
 \beta\in\widehat F,\quad\alpha\in\widehat F\cdot\{1,J\}.
\]
For the order-two factor, reduction of a conjugator in $S_3$ lies in the centralizer of its transposition and is removed by a power of $s$. For the order-three factor, the two possible conjugacy classes are represented by $r,r^{-1}$; their normalizer in the extended group adds $sJ$, which reverses $r$. Comparing reductions in $S_3$ then places the conjugator in the displayed coset union. The determinant of its matrix representative is a single sign $\delta$. This step simultaneously excludes the spurious scalar branch at every prime.

Choose matrix representatives for $\alpha,\beta,f_G$, with determinants $\delta,1,1$, and define the lifts of $\xi,\eta$ by their matrix coboundaries. They are independent of the chosen global signs of $\alpha,\beta$. With
\[
 \xi=\begin{pmatrix}A&B\\C&D\end{pmatrix},\qquad
 \eta=\begin{pmatrix}P&Q\\Q&T\end{pmatrix},
\]
conjugacy of $\xi R$ to $R$ gives $C+D-B=1$; conjugacy of $\eta S$ to $S$ gives the symmetry of $\eta$. The projective loop relation lifts as
\[
 \xi=\varepsilon\eta
 \begin{pmatrix}1&0\\\mu-1&1\end{pmatrix},
 \qquad\varepsilon\in\{\pm1\}\text{ a single global sign}.
\]
Comparing entries gives $B=\varepsilon Q$, $D=\varepsilon T$, and
$1=C+D-B=\mu D$. In particular $D=\mu^{-1}$ is a unit. Put
\[
 V=\begin{pmatrix}\mu&\mu B\\0&1\end{pmatrix}.
\]
A direct multiplication, using $AD-BC=PT-Q^2=1$, proves
\[
 VRV^{-1}=\xi R=\alpha R\alpha^{-1},\qquad
 VSV^{-1}=\varepsilon\eta S=\varepsilon\beta S\beta^{-1}.
\]
The sign in the second equality is essential in the linear group and disappears in its stated projective quotient. Thus $V$ implements $q\circ u_0$ on the two generators $s,r$. For $u=\operatorname{Int}(f_G^{-1})u_0$, take $U=f_G^{-1}V$. Agreement on these topological generators gives agreement of the continuous homomorphisms on all of $\widehat B$. Also $\det U=\mu$.

This formula has the claimed ambiguity. Gauge by $\ell_0^n$ replaces $f_G$ and $V$ by their common left translate, so leaves $U$ unchanged; a global change of sign of $f_G$ changes $U$ by $-I$. For composition, the normalized representatives form the subgroup normalizing $\langle y\rangle$, and their matrices $V$ have lower-right entry one. The product has the same property and implements the composite on $R,S$, so it is the normalized matrix of that composite. Indeed a matrix commuting with both $R,S$ is scalar, as one checks entry by entry, and the lower-right entry fixes the scalar. For general representatives the cocycle formula $f_G''=f_G' u'(f_G)$ supplies precisely the corresponding multiplication law modulo the global sign. Hence $u\mapsto U(u)$ is a homomorphism.

At a finite level, an element of the projective congruence kernel is sent to the identity by the right-hand side, proving invariance. This is a genuine descent proof: the matrix $U$ was constructed before descent was asserted.
\end{proof}

\begin{claim}{C.d2}{Proved}{The norm equations and a correction visible on complex conjugation}
Writing
$\alpha=\left(\begin{smallmatrix}a&b\\c&d\end{smallmatrix}\right)$ and
$\beta=\left(\begin{smallmatrix}r&s\\t&v\end{smallmatrix}\right)$, the matrix loop relation implies
\[
 \delta(c^2+cd+d^2)=\varepsilon(t^2+v^2)=\mu^{-1},
 \qquad
 \delta(ac+bc+bd)=\varepsilon(rt+sv).
\]
It follows algebraically that $\delta\equiv\mu\pmod3$ and $\varepsilon\equiv\mu\pmod4$. These are therefore already constraints on the whole actual profinite solution group of \S46. Scan198 (73) omits a factor $\delta$ in its first equality; the correct later formula is (88), scans202--203.
\end{claim}
\begin{proof}
Multiplying the matrix coboundaries gives
\[
 D=\delta(c^2+cd+d^2),\quad
 B=\delta(ac+bc+bd),\quad
 T=t^2+v^2,\quad Q=rt+sv.
\]
Substitute $D=\mu^{-1}$, $D=\varepsilon T$ and $B=\varepsilon Q$ from C.d1. Modulo three the first norm is $(c-d)^2$, which is one if it is a unit. Hence $\delta\equiv\mu$ there. The second norm is an odd unit; precisely one of $t,v$ is odd, so its residue modulo four is one, giving $\varepsilon\equiv\mu$. No class field theory is needed for these algebraic implications; identifying the signs as Galois characters is a separate arithmetic comparison.

The error in (73) is not merely possible: use complex conjugation with $\mu=-1$, $\alpha=J$, $\beta=f_G=I$. Then $\delta=\varepsilon=-1$ and both norms are one. The printed equality $c^2+cd+d^2=\varepsilon(t^2+v^2)$ reads $1=-1$, already false modulo three. The corrected equality is $c^2+cd+d^2=\delta\varepsilon(t^2+v^2)$. The attached matrix check includes this counterexample. This does not affect (88).
\end{proof}

\begin{claim}{C.d3}{Proved}{The adelic solution locus is a product with global sign compatibility}
The matrix systems of scans204--205 have the following exact coordinate description. For each prime $p$ and fixed global signs $\delta,\varepsilon$, let $X_p^{\delta,\varepsilon}$ consist of
\[
 (\mu_p,\alpha_p,\beta_p,f_p),\quad
 \mu_p\in\mathbb Z_p^\times,\quad
 \det\alpha_p=\delta,\quad\det\beta_p=\det f_p=1,
\]
satisfying
\[
 \alpha_pR\alpha_p^{-1}R^{-1}
 =\varepsilon\beta_pS\beta_p^{-1}S^{-1}
 \begin{pmatrix}1&0\\\mu_p-1&1\end{pmatrix}.
\]
At $p=2$ additionally require $\alpha_2,\beta_2,f_2\equiv I\pmod2$. The unquotiented adelic locus is
\[
 X=\coprod_{\delta,\varepsilon\in\{\pm1\}}
 \left\{(x_p)_p\in\prod_pX_p^{\delta,\varepsilon}:
 \mu_3\equiv\delta\pmod3,\ \mu_2\equiv\varepsilon\pmod4\right\}.
\]
The parameter space of triples as used in the manuscript is its quotient by the simultaneous left action
$(\alpha,\beta,f)\mapsto(L_0^n\alpha,L_0^n\beta,L_0^nf)$, $n\in\widehat{\mathbb Z}$, and the three independent \emph{global} signs on $\alpha,\beta,f$. Local independent signs must not silently replace those three global signs. Its matrix quotient is obtained by
\[
 (\mu,\alpha,\beta,f)\longmapsto
 \left[f^{-1}\begin{pmatrix}\mu&\mu B\\0&1\end{pmatrix}\right],
 \qquad B=(\alpha R\alpha^{-1}R^{-1})_{12}.
\]
This gives a concrete product-of-local-equations description, but does not assert that every point of this locus comes from an actual profinite automorphism.
\end{claim}
\begin{proof}
A matrix over $\widehat{\mathbb Z}=\prod_p\mathbb Z_p$ is invertible precisely when every local determinant is a unit. All the displayed equations use only addition, multiplication and inversion of unit determinants, so they hold if and only if they hold componentwise. C.d2 supplies the two residual sign compatibilities. The full sign $\delta$ must be one of the two elements of the diagonal $\{\pm1\}$ because it comes from the single coset $\widehat F\cdot\{1,J\}$; the same applies to the projective equality sign $\varepsilon$. The simultaneous gauge is coordinatewise because $\widehat{\mathbb Z}$ itself is the unrestricted product. The matrix quotient is gauge invariant by the calculation in C.d1 and is unchanged by the allowed signs projectively.

These equations are the $\widehat{\mathbb Z}$-points of the corresponding signed affine matrix schemes, localized at the indicated determinants. Their coboundary maps are the maps $F_R(\alpha)=\alpha R\alpha^{-1}R^{-1}$ and $F_S(\beta)=\beta S\beta^{-1}S^{-1}$ studied in \S48, scans256--266. Their fibers are right torsors under the centralizers $T_R,T_S$: if two matrices have the same coboundary, their quotient centralizes the indicated generator, and conversely. This proves the direct relationship to the group-scheme calculation, without equating a profinite image with every adelic point. At three $T_R$ and at two $T_S$ have non-torus fibers, so the characteristic-zero description as maximal tori cannot be substituted integrally.
\end{proof}

\begin{claim}{C.d4}{Proved}{The adelic group is an explicit fiber product of normalizers}
Put $K=SJ$, so that $K^2=I$ and $KRK^{-1}=R^{-1}$. For $\mu\in\widehat{\mathbb Z}^{\times}$ write $\delta=\varepsilon_3(\mu)$, $\varepsilon=\varepsilon_2(\mu)$, and $a=(1-\delta)/2\in\{0,1\}$. The gauge-and-sign quotient in C.d3 is canonically the following group, with componentwise multiplication:
\[
 \mathcal G_{\rm ad}=\left\{([T],[W],[U]):
 \begin{array}{l}
 T,W,U\in\operatorname{GL}_2(\widehat{\mathbb Z}),\quad
 \det T=\det W=\det U=\mu,\\
 TRT^{-1}=R^\delta,\qquad WSW^{-1}=\varepsilon S,\\
 T\bmod2=S^a,\qquad W\bmod2=U\bmod2=I
 \end{array}\right\}.
\]
Each bracket denotes quotient by its own single global sign. Thus this is a closed subgroup of a product of two specified matrix normalizers and $\operatorname{GL}_2(\widehat{\mathbb Z})$, modulo three central signs. Equivalently it is the subgroup of the product of its local normalizer systems with common local determinants and the two global sign characters. The isomorphism from the manuscript coordinates is
\[
 A=\alpha K^a,\qquad
 T=A^{-1}V,\quad W=\beta^{-1}V,\quad U=f^{-1}V,
 \qquad V=\begin{pmatrix}\mu&\mu B\\0&1\end{pmatrix}.
\]
The canonical invariants of actual $\widetilde M^!$ therefore define a continuous homomorphism to $\mathcal G_{\rm ad}$. The correction of the $J$-coset multiplication, omitted in scan193, is
\[
 \alpha''=V'\alpha K^a(V')^{-1}\alpha'K^a,\qquad
 \beta''=V'\beta(V')^{-1}\beta',\qquad
 f''=f'U'f(U')^{-1},\qquad\mu''=\mu'\mu.
\]
These are identities modulo the recorded global signs and gauge; the exponent $a$ belongs to the right factor $u$ in $u''=u'u$.
\end{claim}
\begin{proof}
First work before the sign quotient. Since $\det A=1$, all three displayed determinants are $\mu$. The identities in C.d1 give
$A R^\delta A^{-1}=VRV^{-1}$ and
$\beta S\beta^{-1}=\varepsilon VSV^{-1}$.
These are exactly the two normalizer equations for $T,W$. Because $\alpha,\beta,f,V$ reduce to $I$ modulo two and $K\bmod2=S$, the asserted parity conditions follow. A simultaneous left gauge multiplies $A,\beta,f,V$ by the same $L_0^n$, leaving $T,W,U$ unchanged. The three global signs act independently on the three resulting factors.

Conversely, given such $(T,W,U)$, choose $V=\operatorname{diag}(\mu,1)$ and set
\[
 \alpha=VT^{-1}K^a,\qquad \beta=VW^{-1},\qquad f=VU^{-1}.
\]
Their determinants are $\delta,1,1$, and all three reduce to $I$ modulo two. The normalizer equations imply
$\alpha R\alpha^{-1}=VRV^{-1}$ and
$\beta S\beta^{-1}=\varepsilon VSV^{-1}$.
Since $V(SR)V^{-1}=(SR)^\mu$, multiplying these equations gives exactly the loop relation of C.d3. Any other normalized upper-triangular choice is $V_b=\left(\begin{smallmatrix}\mu&b\\0&1\end{smallmatrix}\right)$ with $b\in2\widehat{\mathbb Z}$; it differs by the unique gauge $L_0^{-b/2}$. This proves bijectivity with precisely the stated quotient, including the prime two.

The target is a group intrinsically. Products of $R$-normalizers have exponents $\delta'\delta$; products of $S$-normalizers have signs $\varepsilon'\varepsilon$. Determinants multiply, and both signs are characters of the determinant. Modulo two, powers of the same $S$ multiply as required. The defining equations are closed, so this is a profinite group. This proves associativity without assuming it for an ad hoc triple law.

To see the actual profinite compatibility, now use the lifts $s,r,J$ in the extended free product and put $\kappa=sJ$, so that $\kappa r\kappa^{-1}=r^{-1}$. Although $\alpha$ may lie outside $\widehat B$, the element $A=\alpha\kappa^a$ lies in $\widehat B$, and
$u_0(r)=A r^\delta A^{-1}$. Consequently
$A''=u_0'(A)A'$ is a well-typed cocycle formula on $\widehat B$. C.d1 permits replacing $u_0'$ by conjugation by $V'$ on its matrix quotient. Thus
$A''=V' A(V')^{-1}A'$; returning to $\alpha$ gives the formula in the statement. The formulas for $\beta,f$ follow from their ordinary cocycles in $\widehat F$. Direct substitution then gives $T''=T'T$, $W''=W'W$ and $U''=U'U$. This also proves continuity. The finite program checks the reconstruction and the corrected product for both sign cosets at primes three and five.

Finally, $N_R^\delta=T_RK^a$ and $N_S^\varepsilon=T_SJ^{(1-\varepsilon)/2}$: multiplication by the indicated fixed normalizer reduces either equation to a centralizer equation. Thus the first two factors are exactly the centralizer schemes of \S48 with their inversion components, with the determinant fiber product and the parity condition. Their bad fibers at three and two remain the fibers described in C.d3, not tori obtained by ignoring those primes.
\end{proof}

\begin{claim}{C.d5}{Open}{The remaining image problem in the two centralizer factors}
C.d4 determines the adelic group and the homomorphism to it, but does not show that the image of $\widetilde M^!$ is all of $\mathcal G_{\rm ad}$. The projection to $U$ is onto the matrices reducing to $I$ modulo two, modulo the global sign: the geometric image supplies every determinant-one such matrix, and the arithmetic subgroup supplies every determinant through the cyclotomic character. The undecided part is simultaneous realization of the two remaining normalizer coordinates over a prescribed $U$.

An exact residual instance is the determinant-one fiber over $U=I$: is every pair of classes $T\in T_R$, $W\in T_S$ with determinant one and reduction $I$ modulo two attained by an actual loop automorphism whose congruence action is trivial? Already its reduction modulo three is not determined by the individual norm equations. C.d4 constructs every such adelic matrix solution, and the finite code verifies its reconstruction and multiplication; none of those computations asserts a lift to an automorphism of the full free profinite group. This is the remaining image statement, not a missing definition of the adelic group.
\end{claim}
% anchor:C.d5
\manuscriptanchor{\S 46; fol401; scans204--205}

\begin{proof}[Established projection and attempted attack]
The geometric image is onto the level-two subgroup of $\operatorname{SL}_2(\widehat{\mathbb Z})/\{\pm I\}$ by the elementary-matrix argument in D.3. Given a target $U$ of determinant $\mu$, choose an arithmetic representative whose cyclotomic character is $\mu$; its matrix has that determinant by C.d1 and the arithmetic comparison. Multiplying by a geometric matrix adjusts it to the specified $U$. This proves the projection claim.

For $\mu=1$, the other two factors are the norm-one centralizers. One can freely choose their local points and use the inverse construction in C.d4, so there is no further matrix obstruction. However, a chosen matrix conjugator need not lift to a conjugator of the prescribed element in $\widehat B$: equality after $q$ does not imply equality before $q$. The kernel of $q$, which Part B proves nontrivial, is precisely the missing lifting information. Establishing compatible profinite conjugators and the generation condition of B.m1 is the unresolved step. This pinpoints the smallest requested congruence level at which this image fiber is not determined here, rather than replacing it by a global surjectivity claim.
\end{proof}


\begin{claim}{C.r4}{Proved}{What the congruence proof decides, and its graded limitation}
The projective congruence-kernel invariance in \S39 and the resulting
square-class action follow for all of $\GTz$ by C.d1, with the
normalizations stated there. Therefore these conditions do not
distinguish Galois from $\GT$ or from $\GTz$ in that projective setting.
The signs in C.d2 are algebraic consequences, not new arithmetic
restrictions. The full adelic conditions do not factor through the
single rational pro-unipotent or pro-$\ell$ quotient used in Part B.
\end{claim}
% anchor:C.r4
\manuscriptanchor{\S 39; fols.298 verso--303; scans21--30; \S 41; fols.315--316 verso; scans54--57; \S 46; fols.398--400; scans198--202}

\begin{proof}
The conjugating matrix constructed in C.d1 has determinant $\mu$.
Changing a genuine representative by an inner element changes it by
the corresponding geometric matrix; changing a scalar changes its
determinant by a square. Hence its projective outer class is obtained
from the square class of $\mu$. The arithmetic subgroup realizes every
$\mu$ by its cyclotomic character, so the projective congruence outer
images of all three groups agree. At individual projective levels this
square-class map may have additional kernel, for example at level four;
the linear-group table C.r1 is not used to assert otherwise.

For the graded claim, a finite quotient of a pro-$\ell$ group must be
an $\ell$-group. The congruence image modulo an odd prime generally has
order divisible by several primes, so there is no factorization through
that completion. A characteristic-zero unipotent group has no
nontrivial finite quotient in the category of unipotent algebraic
groups. Individual abelian $\ell$-primary reductions can of course be
read from suitable integral quotients; for instance the level-four
Legendre image is a $2$-group. This does not provide the missing
factorization for the full adelic representation. After imposing
$\lambda=1$ and tensoring the weight-graded Lie algebra with
$\mathbb Q_\ell$, the finite signs have zero differential. We classify
the full congruence constraints as having no faithful graded form in
the specified construction, and the sign differentials as vacuous.
\end{proof}


\begin{claim}{C.r5}{Proved}{The full modular action retains the group-comparison problem}
The cusp-preserving construction of C.a15 induces an injective continuous
homomorphism
\[
 C=\GTz\hookrightarrow\Out(H),\qquad
 H=\widehat{\SL_2(\mathbb Z)}.
\]
Under this injection the arithmetic, pentagon, and two-cycle images are
the respective embedded subgroups $G_\mathbb Q\subseteq\GT\subseteq\GTz$.
Thus their equality on every congruence outer quotient, proved in C.a16,
does not decide their equality on the full modular group. This gives
the precise infinite-group interpretation of C(a), beyond the finite table.
\end{claim}
% anchor:C.r5
\manuscriptanchor{\S \S 45--46; fols.384--391; scans170--184}

\begin{proof}
Suppose a class in $C$ gives an inner automorphism of $H$. Quotienting
by its center gives an inner automorphism of $\widehat B$ with a
conjugator $b\in\widehat B$. On $\widehat F$ its outer class is
the class of $b$ in $\widehat B/\widehat F=S_3$. But the original
class belongs to $C$, the label-preserving centralizer, and the
geometric $S_3$ acts faithfully on the three labels. Their intersection
is the identity. Hence the original class is the identity in
$\Out(\widehat F)$ and in $C$, proving injectivity. Continuity follows
from the continuous finite-quotient constructions and uniqueness of the
central lift; a continuous injection from this compact group has closed
image. The restriction statements follow from the defining inclusions
and arithmetic compatibility in C.a16. This proof identifies the
images as subgroups without claiming a decision on their strictness.
\end{proof}
